\section{Introduction}
\label{sec:intro}

Locating source code that provides concrete evidence for a technical requirement is fundamental in software analysis and capability assessment. 
While classical requirements traceability links formal specifications to code artifacts~\cite{antoniol2002traceability,gao2024triad}, modern software engineering frequently demands verifying whether a codebase implements higher-level framework capabilities (e.g., reactive state management or token authentication). 
We formalize this task as \textbf{Requirement \& Technical Capability Evidence Retrieval}: given natural-language technical requirements---instantiated via mandatory skill specifications from job descriptions (JDs)---the objective is to retrieve the specific code chunks that directly evidence their concrete implementation across multi-file repositories, without requiring explicit requirement IDs.

In retrieval-augmented generation (RAG) pipelines for software engineering, source code must be segmented into chunks for neural encoding. 
Chunking via abstract syntax trees (ASTs) is widely advocated to preserve semantic units (methods, declarations)~\cite{zhang2025cast}. 
However, recent empirical studies on code completion indicate that structure-aware chunking does not reliably outperform simple sliding windows~\cite{wu2026chunking}, and line-based segmentations match syntax-aware splitting across compute budgets~\cite{galimzyanov2025practical}. 
Crucially, prior work conflates chunk boundaries with the contextual metadata (file path, enclosing class, method signature) prepended to chunks. 
Whether structural chunking provides distinct advantages once metadata is controlled remains unexamined for requirement evidence retrieval.

In this Early Research Achievements (ERA) paper, we present a controlled $2 \times 2$ factorial evaluation crossing chunking strategy (node-level AST vs. 50-line windows) with context headers (present vs. absent). We address four research questions:
\textbf{RQ1:} Does prepending a context header improve evidence retrieval from code chunks?
\textbf{RQ2:} Does node-level AST chunking outperform line windows, with and without a context header?
\textbf{RQ3 (Roadmap):} Which header components (file path, class, symbol, call-graph or import context) drive retrieval gains, and how do the trade-offs generalize across languages?
\textbf{RQ4 (Roadmap):} Does retrieved evidence improve human accuracy and confidence in requirement verification?

We evaluate RQ1 and RQ2 across 25 JDs and 40 open-source repositories (20 Java Spring Boot, 20 React/TypeScript; 434 line-window and 438 AST chunks) at three tiers: BM25, dense bi-encoder (BGE-M3), and cross-encoder reranking (BGE-Reranker-Base), evaluated against 882 blinded human-labeled pairs ($\kappa_w = 0.9822$). 
Findings reveal: (1)~headers significantly improve AST chunks ($\Delta = +0.146$, $p_{\mathrm{Holm}} < 0.001$), but not line windows ($\Delta = +0.041$, $p_{\mathrm{Holm}} = 0.2411$); (2)~without headers, line windows significantly outperform AST chunks ($\Delta = -0.239$, $p_{\mathrm{Holm}} < 0.001$); (3)~headerless line windows significantly outperform AST chunks with headers ($\Delta = -0.093$, $p_{\mathrm{Holm}} = 0.0438$); (4)~line windows with headers achieve highest overall effectiveness ($\text{NDCG@10} = 0.621$), outperforming AST with headers ($\Delta = -0.134$, unadjusted $p = 0.0003$); and (5)~with the final parser configuration, the AST fallback rate is $0.0\%$, showing that the performance gap is not caused by fallback degeneration.

\textbf{Contributions:} (1)~\textit{Factorial Framework:} Isolating chunking boundaries from context headers in requirement evidence retrieval. 
(2)~\textit{Empirical Findings:} Evidence that sliding windows match or exceed node-level AST chunking, with headers substantially improving AST chunks (the underlying mechanism is untested; Section~\ref{sec:discussion}). 
(3)~\textit{Curated Benchmark:} A replication package with 882 blinded dual-annotated pairs ($\kappa_w = 0.9822$, 100\% Judged@5, zero orphaned pairs). 
(4)~\textit{ERA Roadmap:} An actionable research agenda addressing header components, dependency graphs, and human practitioner verification.
